\begin{center}
  \textbf{\LARGE Summary Sheet}\\[0.35em]
  \textbf{Team \#M2026232}
\end{center}

\setlength{\parskip}{0.35em}
\setlength{\parindent}{0pt}
\setlist[itemize]{itemsep=0.15em, topsep=0.2em, parsep=0em, partopsep=0em}

\small
\noindent\textbf{Problem.}
Using the provided elevator operational logs, design a practical control policy to:
(i) forecast short-horizon demand, (ii) identify the current traffic regime, and
(iii) proactively park \emph{idle} elevators to reduce passenger waiting and long-wait events.

\vspace{0.55em}
\noindent\textbf{Modeling framework (closed loop).}
We implement a deployable supervisory layer that acts only on idle cars:
\begin{itemize}[leftmargin=1.2em,itemsep=0.25em,topsep=0.25em]
  \item \textbf{Task 1 (Forecasting):}
  Aggregate hall calls into 5-minute slices and predict next-slice demand with a gradient-boosted decision tree (Route A).
  As a transparent benchmark (Route B), we use a time-of-day baseline with regime $AR(1)$ correction.
  \item \textbf{Task 2 (Traffic regime):}
  Discover operational modes from 5-minute feature vectors (intensity, directionality, dispersion, inter-floor tendency, and service pressure)
  using KMeans~\cite{macqueen1967kmeans}, then map clusters to manager-friendly labels (e.g., \emph{Up-Peak}, \emph{Down-Peak}, \emph{Idle/Low}).
  As an interpretable fallback (Route B), we also implement a rule-based regime classifier using the same 5-minute intensity and direction features, providing a transparent baseline for state labeling.
  \item \textbf{Task 3 (Dynamic parking):}
  Learn a mode-conditioned floor-demand distribution and select parking floors for idle elevators by solving a weighted 1D $k$-median problem (Route A).
  As a transparent fallback (Route B), we additionally use a state-aware zone allocator (Lobby/Mid/Upper) with rate-limited, minimal-move repositioning.
  Repositioning is \emph{idle-only} to avoid disrupting in-service dispatch.
\end{itemize}

\vspace{0.55em}
\noindent\textbf{Key quantitative highlight (Task 1 benchmark).}
On the held-out test set, the Route B baseline+$AR(1)$ improves one-step 5-minute demand prediction over the pure time-of-day baseline
(MAE $4.675$ vs.\ $4.781$, RMSE $8.325$ vs.\ $8.528$; about $2.2\%$ and $2.4\%$ relative reductions, respectively).

\vspace{0.55em}
\noindent\textbf{Evaluation summary (Task 3).}
Using a call-level, parameter-consistent simulator, we compare mode-aware dynamic parking against static baselines and a zone-based fallback.
Across overall and operationally critical windows (peak demand, regime transitions, and high-load intervals), the dynamic policy
consistently reduces average waiting time and tail risk (P95/P99) while exposing repositioning cost through an explicit empty-travel metric and using a rate-limited, minimal-move fallback policy for safe deployment.
We further include stress tests to probe failure modes under demand shocks and regime confusion, and we recommend an explicit rollback switch
to the interpretable Route~B controller when mode confidence degrades. Because the long-wait rate is threshold- and parameter-dependent, it can be near zero under the nominal setting; we therefore emphasize tail-aware metrics (P95/P99) and stress tests where tail risk re-emerges under heavy-load and perturbed parameters (Appendix~\ref{app:extra}).




\vspace{0.55em}
\noindent\textbf{Strengths and implementation readiness.}
The approach is (i) data-driven yet interpretable, (ii) lightweight (5-minute decision epochs),
(iii) deployable as a supervisory layer on top of standard dispatch logic, and
(iv) robust through conservative design choices (idle-only repositioning, fallback policy).
For deployment risk control, we monitor signal availability and mode-confidence; when these degrade, the supervisor automatically rolls back to the transparent Route~B controller (Appendix~\ref{app:routeb}) without interrupting in-service dispatch.

\vspace{0.55em}
\noindent\textbf{What is novel here?}
\begin{itemize}[leftmargin=1.2em,itemsep=0.25em,topsep=0.25em]
  \item Mode-aware parking: parking targets adapt to detected regime shifts rather than fixed rules.
  \item Optimization view: parking floors are selected by a demand-weighted 1D $k$-median objective (fast and interpretable).
  \item Engineering safeguards: idle-only moves, rate limits, and a transparent fallback policy.
  \item Fail-safe switching: a simple confidence/stability check to toggle between Route A and Route B for safe deployment.
\end{itemize}


\vspace{0.55em}
\noindent\textbf{Limitations.}
Performance depends on calibration of travel/door-time parameters and on the representativeness of the observed log period.
Field deployment should validate under seasonal and event-driven shifts and monitor repositioning cost (empty travel).